\section{Approximation with exact discrete feasibility}
\label{sec:approximation}
Relative information control has two algorithmic consequences. A linear
criterion, weighted trace, can be optimized on the local graph without fixing
the parameter dimension (Section~\ref{subsec:trace-scheme}). When that
dimension is fixed, several nonlinear criteria can instead share one finite
set of feasible designs (Sections~\ref{subsec:spectral-set}
to~\ref{subsec:matroid-cover}). The two results use different computational
models.

\subsection{A rational weighted-trace scheme}
\label{subsec:trace-scheme}
All inputs in this subsection are rational and explicitly encoded in binary,
including the sensitivity matrices, $J_0\succeq0$, and a common weight
$W\succeq0$. We maximize
\[
 \mathcal I(S)=\tr(WJ(S)).
\]
This is a weighted trace of information, not conventional A-optimality.
An FPTAS is an algorithm, polynomial in input length and $1/\varepsilon$,
that returns a feasible value at least $1-\varepsilon$ times the maximum.

\begin{theorem}[Weighted trace with input-sized dimensions]
\label{thm:trace-fptas}
Suppose the covariance belongs to one of the following promised classes:
\begin{enumerate}
\item The scalar model \eqref{eq:scalar-promises}, with fixed rational upper
bounds $0\le\rho\le\rho_0<1$ and $\overline P/r_{\min}\le B_0$.
\item The full-packet model of Theorem~\ref{thm:block-locality}, supplied in
rational coordinates with
$P_t\preceq B_0V_t$ and
$A_tV_{t-1}A_t^\top\preceq\rho_0^2V_t$, for fixed rational
$B_0\ge0$, $0\le\rho_0<1$.
\item The general-decay model \eqref{eq:decay-promises}, or its supplied
metric version \eqref{eq:metric-decay-promises}, with fixed rational
upper bounds on $C/m$ and $\rho<1$.
\item The partial-packet model \eqref{eq:intrinsic-partial-promises}, with
fixed rational upper bounds $s\le s_0$ and $\gamma\le\gamma_0<1$.
\end{enumerate}
Packet, latent-state, and parameter dimensions may be part of the input.
For rational $0<\varepsilon<1$, an exact cardinality $k$, mandatory and forbidden times, and an optional
integer minimum calendar gap $g\ge1$, there is a deterministic rational
FPTAS, and it detects infeasibility exactly. The polynomial depends on the
fixed promise constants, but its exponent does not grow with these dimensions.
\end{theorem}
\begin{proof}
Each class has a rational envelope
\begin{equation}
 \delta_L\le C_0\beta^{L+1},\qquad 0<\beta<1,
 \label{eq:complexity-envelope}
\end{equation}
with $C_0,\beta$ bounded by fixed constants. For the scalar and full-packet
classes one may use $C_0=2B_0/(1-\rho_0)^2$, $\beta=\rho_0$, from the
weaker gain bound \eqref{eq:gain-delta}; the stronger bound is not needed
for complexity. For general decay use $K_{\mathrm{dec}},\vartheta$ in
\eqref{eq:decay-constants}, formed with the fixed upper promises in place of
$C/m$ and $\rho$ (equivalently, apply the theorem with those larger
decay/floor-ratio bounds). For partial packets put
$\kappa_0=s_0/(1+s_0)$ and use
\[
 C_0=2\kappa_0\left\{\frac1{1-\gamma_0}
       +\frac{s_0\gamma_0}{(1-\gamma_0)(1-\gamma_0^2)}\right\},
 \qquad\beta=\gamma_0.
\]
This follows from $\sqrt{\kappa_0s_0}\le s_0$ in
\eqref{eq:partial-normalized-delta}. Zero signal variance or zero contraction
makes the relevant conditional model exact with $L=0$; no logarithm or
expression involving $1/\log(1/\beta)$ is needed in that case.

For $n\ge1$ choose the first $L\ge0$ for which
$C_0\beta^{L+1}\le\varepsilon/2$, stopping at $L=n-1$ if necessary.
Use rational powers and comparisons. At the cap all histories are complete
and their error is zero. Compute \eqref{eq:local-conditionals} in the
original rational coordinates and assign arc score
$\tr(W W_{t,H})\ge0$ to \eqref{eq:calendar-arc}. Add
$\tr(WJ_0)$ once. Exact longest-path optimization then maximizes
$\mathcal I_L(S)=\tr(WJ_L(S))$ over the graph.

The following polynomial extension of the graph enforces feasibility.
Clamp $g$ to $\min(g,n+1)$, which leaves every feasible schedule unchanged.
Before time $t$ record a counter $c\in\{0,\ldots,g-1\}$ for the number of
still-forbidden times. A skip replaces $c$ by $\max(c-1,0)$; a choose is
permitted only at $c=0$ and sets the next counter to $g-1$. The count starts
at zero, increases by one on a choose, and must end at $k$. Initialize the
counter and history at zero. Remove skips at mandatory times and chooses
at forbidden times. The graph has
$O(n(k+1)g2^L)$ states and arcs and represents exactly the required family,
including $L=0$ and $g>L$. Reachability detects conflicting restrictions.
If $k\notin\{0,\ldots,n\}$ the family is infeasible. For $k=0$, or $n=0$,
return the empty schedule only when all restrictions permit it.

Let $S^*$ maximize the true criterion and $\widehat S$ maximize its
surrogate. Lemma~\ref{lem:residual-transfer} and $W,J_0\succeq0$ imply
\begin{equation}
 \mathcal I(\widehat S)\ge
 \frac{\mathcal I_L(\widehat S)}{1+\delta_L}\ge
 \frac{1-\delta_L}{1+\delta_L}\mathcal I(S^*)
 \ge(1-\varepsilon)\mathcal I(S^*).
 \label{eq:trace-approximation-ratio}
\end{equation}
No division by the optimum is used, so a zero optimum is included.

For the size bound, put $a=\log2/\log(1/\beta)$. Minimality before the
full-history cap gives
\[
 2^L\le\max\{1,(2C_0/\varepsilon)^a\}.
\]
If the cap is encountered earlier this bound remains valid: when $L>0$ the
previous window failed the accuracy test. Thus the graph has polynomial
size. A simple algorithm first forms the whole covariance by rational state
recursions, or reads the supplied dense covariance. A local solve uses at most
the scalar coordinates of $L+1$ packets; the number of histories is
$O(n2^L)$. These matrix operations are polynomial in the explicitly supplied
dimensions. In the metric cases whitening is used only in the proof; no
irrational square root is computed by the algorithm.

It remains to bound bit complexity, not only arithmetic counts.
State recursions expand into sums of products along at most $n$ transitions.
A common denominator formed from input denominators to polynomial powers has
polynomial bit length; the logarithm of the number of expanded coordinate
paths is also polynomial. Hence the covariance entries have polynomial bit
length. Rational principal solves have polynomial-size numerators and
denominators by determinant bounds and can be computed by fraction-free
elimination. Each path value sums at most $n$ rational arc scores and one
prior score. Its denominator bit length is bounded by the sum of those
polynomial bit lengths. Rational comparison is therefore polynomial even
when many paths merge into a state. This proves the Turing bound.
\end{proof}

The theorem takes model sensitivities as input; it does not assert that an
arbitrary nonlinear differential-equation model has polynomial-time sensitivity
evaluation. Promise verification is rational linear algebra: for example,
$B B^\top\preceq a^2I$ tests $\norm{B}_2\le a$ without computing eigenvalues.
For supplied block metrics use the corresponding Schur-complement matrix
inequality. Such tests certify the promises of one finite input; the
fixed-constant assumption restricts a sequence of instances. If contraction
approaches one under grid refinement, the exponent $a$ is no longer uniform.
Arbitrary binary resource budgets need not admit a polynomial-size graph.

\subsection{Scope, hardness, and the scalar predecessor}
\label{subsec:approximation-scope}
Theorem~\ref{thm:trace-fptas} selects complete packets, and this restriction
is substantive. The following conditioning calculation shows why; we do not
claim to introduce hardness of correlated selection.
\begin{proposition}[Individual-channel boundary]
\label{prop:channel-hardness}
Unless $\mathrm P=\mathrm{NP}$, there is no deterministic FPTAS for scalar
information under exact individual-channel cardinality, even at a single time
with covariance condition number at most $5/3$, independent isotropic
observation noise, and the fixed full-packet variance promise $B_0=1$.
\end{proposition}
\begin{proof}
Independent Set is NP-hard on cubic graphs
\citep[Theorem 4.1(a)]{Mohar2001}. For a graph on $m$ vertices with adjacency
matrix $A$, take $R=I+A/12$, scalar sensitivity $\mathbf1$, latent covariance
$P=I/3+A/12$, and observation noise $V=2I/3$. Since
$\norm A_2\le3$, $P\succ0$, $P\preceq(7/8)V$, and
$3I/4\preceq R\preceq5I/4$. The transition promise is vacuous.
For a $k$-subset $S$ let $x=R_{SS}^{-1}\mathbf1$. A Neumann-series bound
from $\norm{A_{SS}/12}_\infty\le1/4$ gives
$\norm x_\infty\le4/3$. Each equation for $x_i$ therefore gives
$x_i\ge1-3(4/3)/12=2/3$. Summing those equations yields
\[
 \mathbf1^\top R_{SS}^{-1}\mathbf1
 =k-\frac1{12}\sum_{i\in S}\deg_S(i)x_i.
\]
This is $k$ for an independent set and at most $k-1/9$ otherwise. With
optional fixed prior $J_0=1$, an FPTAS at
$\varepsilon=1/[18(m+1)]$ must return an independent $k$-set whenever one
exists, since its permitted loss is at most $1/18$. Checking that returned
set would decide the NP-hard problem in polynomial time. Trivial impossible
cardinalities can be checked first. The reduction chooses separate channels
within one packet, which Theorem~\ref{thm:trace-fptas} does not allow.
\end{proof}

For scalar information, strong positive results already exist. The
bounded-treewidth Gaussian message scheme of
\citet[Section 3.2, Theorem 43]{Mahalanabis2012} can be specialized to one
unobservable target and designated observable candidates.
Appendix~\ref{app:scalar-prior} supplies the recurrence adaptation, a rational
normalization and regularization, and the conversion from posterior variance
to information. This yields the scalar approximation guarantee in the
predecessor's spectral-rounding model, so we do not claim a first scalar
FPTAS. The explicit rational finite-history realization in
Theorem~\ref{thm:trace-fptas} avoids importing that numerical model.

For a growing latent block or a growing-rank target, the corresponding
Gaussian separator contains growing-dimensional precision messages, so the
bounded-width runtime does not by itself give the dimension dependence of
Theorem~\ref{thm:trace-fptas}. Nor can independent scalar target problems
be optimized separately when they must share a schedule. This is a scope
distinction, not a lower bound on every older method. Weighted
information and common-support multiple-response approximation are established:
\citet{Liu2016} treats correlated Fisher criteria, while
\citet[Section 5]{Liberty2017} gives sparse multiple-regression guarantees
with enlarged support. The exact cardinality in our statement matters.

\subsection{A relative matrix cover of feasible paths}
\label{subsec:spectral-set}
We now fix the information dimension $p\ge1$. Let $G=(V,E)$ be an explicitly
encoded directed acyclic graph, with source $s$, sink $t$, rational
$Q_e\succeq0$ on each edge and rational $Q_0\succeq0$. For a path $P$ write
$J(P)=Q_0+\sum_{e\in P}Q_e$. This notation is local to the additive graph;
$Q_e$ is an information atom, not an approximate observation precision.

\begin{theorem}[Feasible spectral approximation set]
\label{thm:dag-spectral-set}
For rational $0<\eta<1$, a deterministic algorithm returns a set $\mathcal C$
of actual $s$--$t$ paths such that for every feasible path $P$ some
$\widehat P\in\mathcal C$ satisfies
\begin{equation}
 (1-\eta)J(P)\preceq J(\widehat P)\preceq(1+\eta)J(P).
 \label{eq:spectral-cover}
\end{equation}
The set size and Turing bit complexity are polynomial in the entire input
length and $1/\eta$ for fixed $p$. No positive lower eigenvalue, positive
definite prior, or condition-number bound is assumed. The two matrices in
\eqref{eq:spectral-cover} have exactly the same kernel. Infeasible graphs
return the empty set; if $s=t$ the sole path is empty and is returned exactly.
\end{theorem}

The following normalization lemma serves both the graph algorithm and the
represented-matroid extension below.
\begin{lemma}[Rational normalization with exact ranges]
\label{lem:spectral-normalization}
Consider $m$ rational PSD atoms of order $p$, a common PSD prior, and feasible
objects selecting each atom at most once. Decompose each matrix into at most
$p$ labeled terms $w_jv_jv_j^\top$, where $w_j>0$ and $v_j$ are rational;
retain the atom or prior owning each label. There are $M\le p(m+1)$ labels.
There is a rational trial for each independent set of $r$ labels,
$1\le r\le p$, with the following properties. Retained matrices have ranges contained in a
common subspace $\mathcal S$, a reversible rectangular congruence to
$r\times r$ matrices whose entries have absolute value at most $4p$, and
at most $r$ forced atom owners. Every accepted object has transformed
information at least $I_r$ and range exactly $\mathcal S$. Every feasible
object of positive information rank survives at least one trial.
The algebraic normalization, filters, and owner requirements have polynomial
bit complexity for fixed $p$; feasibility is handled by the later
family-specific algorithms.
\end{lemma}
\begin{proof}
Rational LDL factorization supplies the decomposition. The reference
implementation inspects the leading diagonal of each residual. If that
entry is zero, its row and column vanish because
$|a_{ij}|^2\le a_{ii}a_{jj}$, and the coordinate is skipped. Otherwise the
pivot is positive and leaves a rational PSD Schur complement. Each step
removes one coordinate, proving termination even for singular matrices,
without square roots. Fix an independent label set with column matrix
$V_b=[v_{b_1},\ldots,v_{b_r}]$, and define
\begin{align}
 L_b&=(V_b^\top V_b)^{-1}V_b^\top,&\Pi&=V_bL_b,\nonumber\\
 T&=\diag(\tau_i)L_b,&K&=V_b\diag(\tau_i^{-1}),
 \qquad 1\le\tau_i^2w_{b_i}<4.
 \label{eq:spectral-normalization}
\end{align}
Here each $\tau_i$ is a power of two, found by rational comparisons. We have
$TK=I_r$, $KT=\Pi$, and $\Pi$ is the orthogonal projector onto
$\mathcal S=\range V_b$. Reject the trial if $\Pi Q_0\ne Q_0$; discard
an atom if $\Pi Q_e\ne Q_e$. For any retained symmetric matrix,
\begin{equation}
 Q=K(TQT^\top)K^\top.
 \label{eq:spectral-reconstruction}
\end{equation}
Indeed $\Pi Q=Q$ implies $Q\Pi=Q$, so the right side is $\Pi Q\Pi$.
Put $A_e=TQ_eT^\top$ and $A_0=TQ_0T^\top$. Reject the trial if some
diagonal of $A_0$ exceeds $4p$, and discard atoms with that property.
Positive semidefiniteness bounds all remaining entries by $4p$ in magnitude.
Require all distinct nonprior owners of the basis labels. Their factors are present with their original multiplicities,
so every accepted object's information $A$ obeys
\begin{equation}
 A\succeq\diag(\tau_i^2w_{b_i})\succeq I_r.
 \label{eq:forced-lower-floor}
\end{equation}
Together with the range tests this gives exact range $\mathcal S$.

To prove coverage, fix an object whose information has rank $r>0$. The
factors in its atoms and prior span its information range: a vector is in
the kernel of a sum of PSD matrices exactly when it is in every summand's
kernel. Introduce $u_j=\sqrt{w_j}v_j$ for this proof only, and choose $r$
such vectors of maximum $r$-dimensional volume. For their matrix $B$, every
other selected factor has coordinates $u=Bx$ with $|x_i|\le1$. Replacing
column $i$ by $u$ multiplies volume by $|x_i|$, which proves the assertion.
The formal proof carries out this determinant replacement in coordinates
of the factor span, including when it is a proper subspace of the ambient
space. The algorithm enumerates this basis's rational labels. Its range tests
preserve the object, and $TB=\diag(\tau_i\sqrt{w_{b_i}})$ has diagonal in
$[1,2)$. Every transformed selected factor therefore has coordinates of
magnitude less than two. Each matrix has at most $p$ factors, so its
transformed diagonal is at most $4p$ and it survives every filter. Its forced
owners were selected by construction.

Exact LDL, rational inverses, projectors and filters have polynomial bit
length by determinant bounds. A very small rational pivot changes the dyadic
exponent by at most its encoding length; the dyadic rational itself still
has polynomial encoding length. Thus no square-root or maximum-volume oracle
is used, and all actual trial operations are polynomial in bit complexity.
\end{proof}

\begin{proof}[Proof of Theorem~\ref{thm:dag-spectral-set}]
After the stated reachability and $s=t$ cases, let $N=\max(1,|V|-1)$ bound
path length. Apply Lemma~\ref{lem:spectral-normalization} with edge owners.
For a rank-$r$ trial use $h=\eta/(rN)$ and the signed integer labels
\begin{equation}
 z_{e,ij}=\left\lfloor A_{e,ij}/h\right\rfloor,
 \qquad 1\le i\le j\le r.
 \label{eq:signed-profile}
\end{equation}
Do not round the prior. In topological order keep one actual path at every
state consisting of a vertex, a bit mask of forced owners already visited,
and the sum of upper-triangular integer labels. Transitions add labels and
update masks. Every continuation depends only on this state. Strict
topological increase prevents a suffix from revisiting either prefix, so
substituting a representative preserves feasible continuations as well as
their labels. At $t$ return
all representatives with complete masks, across all trials.

Fix a feasible rank-$r$ path and its surviving trial from the lemma. Its
terminal state has a representative $\widehat P$. For each entry a path's
rounding residual lies in $[0,Nh)$; this holds for negative entries too
because the operation is floor. Equal integer sums and the common prior give
$|A(\widehat P)_{ij}-A(P)_{ij}|<Nh$. The spectral norm of their symmetric
difference is at most the maximum absolute row sum, hence at most
$rNh=\eta$. Since $A(P)\succeq I_r$, this proves the two-sided relative
inequality in normalized coordinates. Congruence by $K$ and
\eqref{eq:spectral-reconstruction} prove \eqref{eq:spectral-cover}.
For rank zero, if $Q_0=0$, search the subgraph of exactly zero-matrix edges
and include one path when possible. A zero PSD sum has only zero summands,
so this covers every zero-information path.

For an explicit size bound, put
$d_r=r(r+1)/2$ and
\[
 C_r=\left\lceil8prN^2/\eta+N\right\rceil+2.
\]
Every signed accumulated coordinate has at most $C_r$ values, because
$|A_{e,ij}|\le4p$ and at most $N$ labels are added. There are at most
$|V|2^rC_r^{d_r}$ states and $|E|2^rC_r^{d_r}$ edge extensions per trial,
up to dictionary and fixed-dimensional arithmetic costs. In particular,
\begin{equation}
 |\mathcal C|\le1+\sum_{r=1}^p\binom Mr\,2^rC_r^{d_r}.
 \label{eq:dag-cover-size}
\end{equation}
Normalization, floors and integer sums have polynomial bit length; stored
path reconstruction has length at most $N$. For the reference list
implementation, write $B_r=2^r C_r^{d_r}$. The retained table has at most
$|V|B_r$ entries, while a transient candidate list can contain
$1+|E|B_r$ paths. A full-scan bound on key comparisons is
$|V|(1+|E|B_r)^2$. Key evaluation, edge scans, rational arithmetic, and
copying whole paths contribute additional polynomial factors; state and
extension counts alone are not total memory or bit-work bounds. This gives
a polynomial realization without assuming an optimized dictionary.
Exact ranges and \eqref{eq:forced-lower-floor} prove the kernel claim.
\end{proof}

Exact ranges cannot be omitted. For $a=(1,0)^\top$ and
$b=(1,t)^\top$ with nonzero rational $t$, the matrices $aa^\top$ and
$bb^\top$ approach each other entrywise as $t\to0$ but have different
kernels. Neither can dominate a positive multiple of the other. Also,
$\mathcal C$ need not be a Pareto-optimal set: deleting a dominated matrix
may destroy the upper half of \eqref{eq:spectral-cover}.
The exponent in \eqref{eq:dag-cover-size} grows with $p$. This is a
complexity theorem, not a practical running-time claim or a fixed-parameter
tractability statement with a dimension-independent exponent.

\subsection{Criteria, singular information, and finite-memory transfer}
\label{subsec:spectral-consequences}
One computed set supports objectives chosen afterwards, provided their
values can be compared. If $\Phi\ge0$ is Loewner-nondecreasing and
homogeneous of degree $a>0$, the best representative has value at least
$(1-\eta)^a\max_P\Phi(J(P))$. This follows by applying $\Phi$ to the lower
sandwich; concavity is unnecessary. Monotonicity alone supplies no such
multiplicative guarantee.

For D-optimality, $\eta=\varepsilon/p$ gives determinant ratio
$(1-\eta)^p\ge1-\varepsilon$ by Bernoulli's inequality. Rational
determinants can be compared exactly. Determinant roots have ratio
$1-\eta$; these are not multiplicative log-determinant statements.
For E-optimality the minimum eigenvalue has ratio $1-\eta$. The reference
comparator computes rational characteristic coefficients of the Kronecker
difference $A\otimes I-I\otimes B$. A coefficient bound separates every
nonzero eigenvalue difference from zero. Exact rational PSD threshold tests
and a prescribed finite dyadic bisection then distinguish strict order and
equality, including repeated eigenvalues. The degree is fixed when $p$ is
fixed; no exact-real comparison oracle is supplied. If all paths are
singular the E optimum is zero.

For conventional A-optimality give singular information infinite cost.
Inversion of the lower sandwich gives
$\tr(J(\widehat P)^{-1})\le(1-\eta)^{-1}\tr(J(P)^{-1})$ on positive
definite matrices. The choice $\eta=\varepsilon/(1+\varepsilon)$ gives
ratio $1+\varepsilon$; inverse traces are rational. Existence of a positive
definite feasible matrix is detected by the set itself. For a specified
contrast $c$, give $c^\top J^\dagger c$ infinite cost unless
$c\in\range J$. The common-kernel property permits inversion on the
common range and gives
\[
 \frac{c^\top J(P)^\dagger c}{1+\eta}
 \le c^\top J(\widehat P)^\dagger c
 \le\frac{c^\top J(P)^\dagger c}{1-\eta}.
\]
This is not an assertion that pseudoinverses reverse order across different
kernels. Multiple fixed estimable contrasts with nonnegative weights follow
by summation. Adding any common $H\succeq0$ preserves
\eqref{eq:spectral-cover}, as does any common congruence. The same set can
therefore be reused after either operation, with the stated objective-evaluation
qualifications.

\begin{corollary}[Cover for the true finite-memory design model]
\label{cor:true-spectral-cover}
Suppose an explicit rational PSD-edge graph satisfies, for every feasible
schedule, $(1-\delta)J(S)\preceq J_L(S)\preceq(1+\delta)J(S)$,
$0\le\delta<1$. The set from Theorem~\ref{thm:dag-spectral-set} satisfies
\begin{equation}
 \frac{(1-\eta)(1-\delta)}{1+\delta}J(S)
 \preceq J(\widehat S)\preceq
 \frac{(1+\eta)(1+\delta)}{1-\delta}J(S)
 \label{eq:true-spectral-cover}
\end{equation}
for a representative of every feasible $S$. Under any fixed-promise class
of Theorem~\ref{thm:trace-fptas}, fixed $p$ gives a deterministic polynomial
relative cover, and hence the stated D, E, A and contrast schemes, for the
true selected covariance.
\end{corollary}
\begin{proof}
Apply the local sandwich once to each of $S,\widehat S$ and the graph
cover between them to obtain \eqref{eq:true-spectral-cover}. Given a desired
$0<\xi<1$, take $\eta=\xi/4$ and a window with $\delta\le\xi/8$.
The lower factor's deficit is at most $\eta+2\delta\le\xi/2$. The upper
factor's excess is $(\eta+2\delta+\eta\delta)/(1-\delta)\le17\xi/28<\xi$.
The graph construction and cover are polynomial under the fixed promises;
the finite-history cap is handled as in Theorem~\ref{thm:trace-fptas}.
Choose $\xi$ according to the desired scalar criterion ratio. All these
sandwiches preserve kernels, including singular priors.
\end{proof}

\paragraph{Lean verification boundary.}
The accompanying Lean formalization covers the explicit-DAG
construction from the original rational prior and atoms on a graph with
verified topological vertex indices: the two-sided cover and exact ranges,
finite output counts, and the criterion and conditional transfer arguments
above. It contains executable rational normalization and a
representative-merging path DP; it is not a theorem conditional on a supplied
successful trial or cover certificate. The complete bit-work bound starts
from the original matrix and accuracy encodings and includes factor
initialization, rejected trials, exact arithmetic, stored-data access,
dictionaries and path recovery. Exact D, E, A and rational contrast selection
also have proved execution bounds. The cost model uses schoolbook arithmetic
and explicit storage and scan charges; it gives no wall-clock bound for Lean's
runtime. The upstream stochastic locality estimates and finite-history graph
construction required by the full Corollary~\ref{cor:true-spectral-cover} are
not part of this verification. The represented-matroid theorem below is
formalized separately.

\subsection{Rationally represented matroids and contribution boundary}
\label{subsec:matroid-cover}
A different feasible family admits the same cover, through an established
exact profile algorithm. Let $A\in\Q^{a\times m}$ represent a matroid:
a set is independent precisely when its columns of $A$ are independent.
Its rank $q=\rank A$ is input-sized. Give its elements rational PSD atoms
$Q_e$ of fixed order $p$, and define $J(B)=Q_0+\sum_{e\in B}Q_e$ for bases.
\begin{theorem}[Represented-matroid spectral cover]
\label{thm:matroid-spectral-set}
For the preceding rational representation, one can deterministically return
actual bases satisfying \eqref{eq:spectral-cover} for every target base,
with polynomial output size and Turing bit complexity in input length and
$1/\eta$ for fixed $p$. Singular information and an input-sized matroid
rank are allowed. The same criterion and reuse consequences apply.
\end{theorem}
The complete proof in Appendix~\ref{app:matroid-proof} combines
Lemma~\ref{lem:spectral-normalization} with the squared-determinant profile
polynomial, interpolation and deletion recovery of \citet{Berstein2008};
those algebraic algorithms are prior work. Our construction uses an extra
owner-count coordinate instead of contraction, cached polynomial-time
determinant evaluation, and tensor Lagrange interpolation. The Lean
formalization proves the assembled construction from the supplied rectangular
rational representation, including row reduction, exact singular ranges,
polynomial output size, polynomial charged bit work at fixed $p$, and exact
criterion selectors with their work bounds. Its cost model uses schoolbook
arithmetic and finite storage and scans, and excludes the construction of
proof and observer transcripts. Rational Vandermonde
representations include uniform matroids and their direct sums (partition
quotas); oriented incidence matrices include graphic matroids and spanning
forests, including labelled parallel edges and loops.
The result requires the supplied characteristic-zero representation. It does
not construct one from an independence oracle or finite-field representation,
and it does not cover matroid intersections. Most importantly, additive
matroid atoms do not encode arbitrary correlated history or an additional
matroid restriction on the calendar graph without a further reduction.

Approximate Pareto sets and additive profile schemes are established
\citep{Papadimitriou2000,Tsaggouris2009,Mittal2013}. Their nonnegative
coordinate formulations do not directly turn relative entrywise errors into
relative PSD errors, because cancellation and singular ranges matter.
\citet{Brown2024} already normalize, filter large vectors, force a guessed
subset, and handle several monotone homogeneous criteria in a randomized
fixed-dimensional PTAS for the broader independence-oracle matroid model.
Our rationally represented specialization has deterministic polynomial
accuracy dependence; it does not subsume their oracle model.
Generic decomposition of multiresponse PSD atoms into tied rank-one copies
is also established \citep[Theorem 1]{Filova2026}; decomposition itself is
not a contribution here.

Covariance-state pruning predates these design schemes
\citep{Vitus2012,Atanasov2014}. Their additive-error results concern dynamic
Riccati trajectories and process-noise/covariance bounds. Multiplicative
relaxed dynamic programming is older still
\citep{Alriksson2005,Lincoln2006}. A quadratic lower-envelope guarantee has
a different quantifier: a retained matrix may depend on the tested direction.
For example,
\[
 \min\{x^\top\diag(0,2)x,x^\top\diag(2,0)x\}\le x^\top Ix
 \quad\text{for every }x,
\]
yet neither retained singular matrix is a relative PSD approximation to $I$.
Objective-independent spectral spanners and coresets are also prior art
\citep{Indyk2020,Mahabadi2026}; they retain elements and use fractional
replacements or criterion-specific rounding, rather than returning one
feasible integral object within $1\pm\eta$ of each target matrix.
The projected-vertex method of \citet{Onn2004} addresses convex
maximization. Its vertices are insufficient here: among the three
feasible matrices $\diag(1,3),2I,\diag(3,1)$, the middle projected profile
uniquely maximizes both determinant and minimum eigenvalue.
Fixed-decision-dimensional polynomial optimization \citep{DeLoera2008}
does not supply attainability of these fixed-dimensional matrix profiles:
the number of path or element decisions is still input-sized. Recent
one-exact Pareto frameworks \citep{Bokler2026} likewise require a suitable
positive-objective restriction oracle; they do not automatically supply
the signed-profile to PSD-order reduction used here.

To the best of our knowledge, the specific combined guarantee in
Theorems~\ref{thm:dag-spectral-set} and \ref{thm:matroid-spectral-set}---a
polynomial set of feasible integral objects giving an all-target, two-sided
relative PSD cover, exact singular ranges, and deterministic polynomial
accuracy dependence under the stated explicit representations---is not
stated in these predecessors. The contribution is the
normalization-to-bounded-profiles reduction and its complete guarantees,
not dynamic programming, profile interpolation, PSD pruning, or objective
independence separately. The large polynomial exponents rule out any
inference of practical superiority; Section~\ref{sec:computation} uses
support pricing and certificates instead of these spectral sets.
